\section{Method}
\textbf{Grasp maps.} As in~\cite{kumra2020antipodal}, the network predicts four $224{\times}224$ maps: quality $Q$, $\cos2\Theta$, $\sin2\Theta$ and opening $W$. Targets fill the central third of each ground-truth rectangle. The grasp is read at the most interior pixel of the near-maximum plateau of the smoothed $Q$, with $\theta=\tfrac12\operatorname{atan2}(\sin2\Theta,\cos2\Theta)$ and $h=w/2$.

\textbf{FiLM.} The encoder maps the image to $F\in\mathbb{R}^{128\times56\times56}$, followed by five residual blocks and a decoder. With $t$ the normalised CLIP text embedding, each residual block $k$ is followed by
\begin{equation}
F \leftarrow \big(1+\Delta\gamma_k(t)\big)\odot F + \beta_k(t),
\end{equation}
where $\Delta\gamma_k,\beta_k\in\mathbb{R}^{128}$ come from a zero-initialised linear layer, so each FiLM layer is the identity at initialisation.

\textbf{Language gate (proposed).} FiLM applies the same channel modulation at every pixel, so it can favour one object among several only indirectly. The language gate therefore compares the $14{\times}14$ post-LayerNorm patch tokens $P$ of a frozen CLIP ViT-B/16 with the prompt and re-weights the encoder output:
\begin{equation}
\begin{gathered}
S_{ij}=\cos(W_p\,p_{ij},\,W_t\,t),\\
A=\mathrm{Up}_{56}\big(\sigma(\tau S)\big),\qquad F'=F\odot(1+A),
\end{gathered}
\end{equation}
where $p_{ij}\in\mathbb{R}^{768}$ is a patch token, $W_p$ and $W_t$ are learned affine maps to $\mathbb{R}^{256}$, $\tau$ is a learned scale (initially 10), and $\mathrm{Up}_{56}$ bilinearly upsamples the $14{\times}14$ map to $56{\times}56$. As $A\in(0,1)$, $1{+}A$ never suppresses features. The \textbf{C0 control} replaces $W_t t$ with a learned vector $q$: its gate sees the same CLIP features but not the prompt (FiLM still receives $t$), and it has 131k fewer parameters than C. \textbf{Loss:} Smooth-$\ell_1$ summed over the four maps. Variant D adds $0.5\,\mathrm{BCE}(\tau S, M)$, $M$ the pooled ground-truth $Q$ mask.
